%==============================================================================
%  PESAD - Prolog Expert System for Anxiety, Obsessive-Compulsive and Trauma-Related Disorders
%  Technical report (English edition)
%==============================================================================
\documentclass[11pt,a4paper,oneside]{report}

\usepackage[utf8]{inputenc}
\usepackage[T1]{fontenc}
\usepackage[english]{babel}
\usepackage{lmodern}
\usepackage[a4paper,margin=2.6cm]{geometry}
\usepackage{microtype}
\usepackage{amsmath,amssymb}
\usepackage{graphicx}
\usepackage{booktabs}
\usepackage{longtable}
\usepackage{array}
\usepackage{enumitem}
\usepackage{xcolor}
\usepackage{listings}
\usepackage{caption}
\usepackage{fancyhdr}
\usepackage[hidelinks,colorlinks=true,linkcolor=darkblue,citecolor=darkblue,urlcolor=darkblue]{hyperref}

\definecolor{darkblue}{RGB}{20,40,110}
\definecolor{codegray}{RGB}{120,120,120}
\definecolor{codekw}{RGB}{160,30,90}
\definecolor{codestr}{RGB}{20,110,60}
\definecolor{codebg}{RGB}{248,248,248}

% ---- Prolog listing style ---------------------------------------------------
\lstdefinelanguage{Prolog}{
  morekeywords={is,mod,rem,div,not,fail,true,false,dynamic,multifile,discontiguous,
    assert,asserta,assertz,retract,retractall,findall,member,length,nth1,
    predsort,min_list,max_list,format,write,writeln,nl,read,repeat,halt,consult,reconsult},
  sensitive=true,
  morecomment=[l]{\%},
  morestring=[b]',
}
\lstset{
  language=Prolog,
  backgroundcolor=\color{codebg},
  basicstyle=\ttfamily\footnotesize,
  keywordstyle=\color{codekw}\bfseries,
  commentstyle=\color{codegray}\itshape,
  stringstyle=\color{codestr},
  numbers=left,
  numberstyle=\tiny\color{codegray},
  numbersep=8pt,
  frame=single,
  rulecolor=\color{codegray!40},
  breaklines=true,
  breakatwhitespace=true,
  showstringspaces=false,
  columns=fullflexible,
  tabsize=2,
  captionpos=b,
  xleftmargin=14pt,
  framexleftmargin=14pt,
}

\pagestyle{fancy}
\fancyhf{}
\fancyhead[L]{\small\itshape PESAD --- Prolog Expert System (DSM-5-TR)}
\fancyhead[R]{\small\thepage}
\renewcommand{\headrulewidth}{0.4pt}
\raggedbottom

\newcommand{\cf}{\textit{CF}}
\newcommand{\code}[1]{\texttt{#1}}

%==============================================================================
\begin{document}

\begin{titlepage}
  \centering
  \vspace*{1.5cm}
  {\scshape\LARGE University of Pisa\par}
  \vspace{0.3cm}
  {\scshape\large Department of Computer Science\par}
  \vspace{0.2cm}
  {\itshape Fundamentals of Artificial Intelligence\par}
  \vfill
  {\Huge\bfseries PESAD\par}
  \vspace{0.4cm}
  {\Large Prolog Expert System for Anxiety, Obsessive-Compulsive\\ and Trauma-Related Disorders\par}
  \vspace{0.8cm}
  {\large A rule-based expert system for the diagnosis of mental\\ disorders under uncertainty (DSM-5-TR)\par}
  \vfill
  {\large\bfseries Donato Meoli\par}
  \vspace{0.5cm}
  {\large Technical report\par}
  \vfill
\end{titlepage}

\pagenumbering{roman}
\begin{abstract}
\noindent
PESAD (\emph{Prolog Expert System for Anxiety, Obsessive-Compulsive and
Trauma-Related Disorders}) is a rule-based expert system, written in SWI-Prolog,
that supports the diagnosis of the disorders of the DSM-5-TR Anxiety,
Obsessive-Compulsive and Related, and Trauma- and Stressor-Related chapters. The
system reasons under uncertainty by
means of \emph{certainty factors}, propagating a degree of belief from the
elementary symptoms reported by the clinician up to the candidate diagnoses. It
offers two reasoning modes ---an \emph{investigation} mode that ranks all
plausible disorders and a \emph{control} mode that verifies a single
hypothesis--- and it is able to justify its conclusions, both by explaining
\emph{why} a question is asked and by displaying the proof tree that shows
\emph{how} a conclusion was reached. This report describes the theoretical
background of expert systems, the analysis of the clinical domain, the
acquisition and conceptualization of the knowledge, and the implementation of
the system: its modular architecture, the knowledge-acquisition protocols, the
certainty-factor calculus, the explanation module and the bilingual user
interface.
\end{abstract}

\tableofcontents
\clearpage
\pagenumbering{arabic}

\chapter{Introduction}

The aim of this work is to present a study concerning the development of an
expert system, built in a logic programming environment, for the diagnosis of
anxiety disorders.

Throughout the report we document the various design choices adopted during the
definition phase, as well as the behaviour that the expert system exhibits
during its execution.

Before discussing the project in detail, we provide a brief introduction to the
features that characterise an expert system, to expert systems applied to the
medical field, and to the handling of uncertain knowledge.


\section{Expert Systems}

One of the most actively developed areas in Artificial Intelligence is the study
and design of expert systems, also referred to as \emph{Expert Systems} or
\emph{Knowledge-Based Systems}. In the following sections we illustrate the
features and the components that characterise an expert system.


\subsection{Definition}

The term \emph{expert system} was introduced in 1977 by Feigenbaum, who also
provided a precise definition:

\begin{quotation}
An expert system is a computer program that uses knowledge and reasoning
techniques to solve problems that would normally require the help of a human
expert. An expert system must be able to justify or explain why a particular
solution was reached for a given problem.
\end{quotation}

An expert system must be capable of emulating the work of a human expert in a
given domain. In particular, an expert system must be able to perform the same
actions, provide the same judgements, and offer the same explanations as a
human. This feature, however, should not obscure the ultimate goal of an expert
system. Indeed, contrary to what one might think, its purpose is not to replace
the expert, but rather to assist and help them improve their performance.

The fundamental features of an expert system are the following:

\begin{itemize}
    \item \emph{Possession of specific knowledge}: it is important that an expert
    system possesses enough domain-specific knowledge. The quality of an expert
    system, in fact, does not depend only on the formalism or the inference
    technique adopted, but also on the knowledge that the program possesses;

    \item \emph{Adherence of the domain to reality}: the domain in which an
    expert system operates must not be a simplified model of the real domain, but
    must represent it as it is, of course within certain limits, yet never
    resorting to excessively abstract models;

    \item \emph{Reasoning capability}: an expert system must draw inferences and
    make decisions by exploiting the available knowledge base;

    \item \emph{Capability of interacting with the external world}: an expert
    system must have a communication protocol with the external world. Such
    communication is necessary in order to retrieve the information required to
    carry on the reasoning and to provide explanations to the user;

    \item \emph{Capability of explaining its own operation}: a good expert system
    must be able to make explicit and explain the decisions taken and the
    reasoning by which it reached a conclusion. This capability strengthens the
    degree of trust that the user places in the system;

    \item \emph{Learning capability}: a feature that, while not common to all
    expert systems, is very interesting to consider. Just as a human expert
    learns new experiences and broadens their knowledge over time, an expert
    system should likewise be able to learn new knowledge.
\end{itemize}


\subsection{Components}

Within an expert system it is possible to identify four fundamental components:
the knowledge base, the inference engine, the user interface, and the working
memory.

\subsubsection*{Knowledge Base}

The knowledge base represents the set of information needed to address and solve
a problem in a specific domain.

Such information is called \emph{knowledge}, and it must be represented by means
of a formalism that the machine can interpret (for example, Horn logic,
description logic, modal logic, and so on).

Typically, the knowledge base is organised into two blocks:

\begin{itemize}
    \item \emph{Block of assertions or facts}: it represents the declarative
    knowledge related to a particular problem to be solved, that is, true facts
    introduced at the beginning of the consultation;

    \item \emph{Block of relations or rules}: it contains the relations among the
    elementary facts that are useful for building the body of knowledge. These
    are usually expressed in the form of production rules, or according to other
    types of representation such as semantic networks, frames, scripts, and so
    on.
\end{itemize}

\subsubsection*{Inference Engine}

The inference engine is a program able to deduce, starting from the knowledge
base, the conclusions that constitute the solution of the problem in the domain
in which the expert system operates.

The inference engine is able to identify the set of rules that are applicable at
a given moment, to select the rule to apply among the candidate ones, and to
execute the action indicated by the selected rule.

\subsubsection*{User Interface}

The user interface is responsible for facilitating the interaction between the
user and the expert system. It can be built using user-friendly graphical
interfaces or command-line interfaces. It is important that the interface makes
the dialogue with the user as close as possible to a natural one.

\subsubsection*{Working Memory}

The working memory is the component that allows the storage of newly observed
evidence or of data useful for solving the problem during a specific working
session. The content of this memory will therefore change throughout the entire
execution session of the expert system.


\subsection{Roles}

In general, the roles that interact with an expert system are three:

\begin{itemize}
    \item \emph{Domain expert}: the actor who holds the knowledge and experience
    about a specific domain and is able to solve the problem at hand;

    \item \emph{Knowledge engineer}: the actor who encodes the knowledge,
    provided by the domain expert, into a formalism that the expert system can
    understand;

    \item \emph{User}: the individual who consults the expert system in order to
    benefit from it.
\end{itemize}


\subsection{Building Environments}

The construction of an expert system can be carried out using two different types
of environment:

\begin{itemize}
    \item \emph{Skeleton environments}: these are environments derived from
    existing expert systems, from which the knowledge related to the application
    domain has been removed. Examples include Emycin, Kas, Expert, and so on;

    \item \emph{Environment-type environments}: these are environments that
    provide a language for representing knowledge. Examples include Clips,
    Prolog, and so on.
\end{itemize}

Comparing the two types of environment, one can state that ``Skeleton'' systems
are more immediate, but less flexible, since they are affected by the
specialised architecture of the original product; ``Environment''-type systems,
on the other hand, have flexibility as their strength, while their weakness is
the greater complexity.


\section{Expert Systems in Medicine}

One of the most successful fields in which expert systems have been applied is
the medical one.

Typically, a medical expert system can be used in two different ways: as
decision support or as decision acquisition.

\begin{itemize}
    \item \emph{Decision Support}: expert systems of this type aim to suggest
    options or diagnoses to an expert who bears the responsibility of making
    decisions;

    \item \emph{Decision Acquisition}: expert systems for decision acquisition
    aim to allow a non-expert person to make a decision that goes beyond their
    level of knowledge and experience.
\end{itemize}

The reasons that have driven the spread of expert systems in the medical field
are essentially to be found in the possibility of incorporating, within a
knowledge base, the entire domain-specific experience, isolating it from
possible errors or oversights that may undermine a diagnosis carried out by the
human expert alone. In an expert system, indeed, if a line of reasoning is
correct, it will remain so across subsequent consultations as well.


\section{Handling Uncertain Knowledge}

The cognitive mechanism of the human expert often relies on empirical rules
acquired from their own experience and on data that are not completely certain.

An expert is able to reason even when the available information is uncertain or
incomplete. This implies that conclusions obtained from knowledge held to be
generally true, and by means of approximate reasoning, are not always absolutely
correct.

Starting from these premises, the need emerges for approaches that are able to
replicate this approximate reasoning within expert systems as well.

There are several ways to represent and formalise uncertainty: they range from
quantitative methods (certainty factors, Bayesian networks, belief functions,
the fuzzy approach) to symbolic methods (non-monotonic logics).

All these methods share the common feature of having to deal with the
combination of evidence and the propagation of uncertainty from facts to
conclusions.

%==============================================================================
%  Chapter: The PESAD System
%  Section: Domain Analysis
%==============================================================================
\chapter{The PESAD System}

The following sections present the analysis of the domain, the phases of
identification, conceptualization and knowledge representation and, finally, the
implementation of the functional components of the realized expert system.

\section{Domain Analysis}

This section discusses the domain under study for the construction of the expert
system. A brief introduction to the concept of ``anxiety'' is followed by its
distinction into ``state anxiety'' and ``trait anxiety''. The symptoms of
anxiety are then introduced, together with their classification according to the
functional system they affect. Finally, the classification adopted in the
present work is described: that of the fifth edition, text revision, of the
\emph{Diagnostic and Statistical Manual of Mental Disorders}
(DSM-5-TR, 2022)~\cite{dsm}, which distributes the conditions of interest across
three distinct chapters, each of which is analysed in turn.

\subsection{Anxiety}

From an etymological point of view, anxiety derives from the Latin term
\emph{anxia} (to tighten). Anxiety is a universal emotion which, in itself, it
would not be inappropriate to experience, since it represents a necessary part
of the response to stress. It constitutes, in fact, a defence mechanism aimed at
anticipating the perception of danger even before it has clearly manifested
itself. Anxiety typically sets in motion the physiological mechanisms that push,
on the one hand, towards exploration in order to identify the danger and confront
it in the most appropriate manner and, on the other hand, towards avoidance and
flight.

Anxiety is generally regarded by mental-health professionals as ``the mother of
all emotions'' and, precisely for this reason, it is useful to differentiate an
anxiety understood as an existential condition (``trait'' anxiety) from a
pathological anxiety (``state'' anxiety).

The distinction between the concepts of ``trait'' anxiety and ``state'' anxiety
was introduced by Cattell and Scheier (1961) and further elaborated by
Spielberger and colleagues in the development of their self-assessment scale, the
State--Trait Anxiety Inventory~\cite{stai}. The two paragraphs below report some
salient points of the two concepts of anxiety, as elaborated by Spielberger and
colleagues.

\subsubsection{Trait Anxiety}

Trait anxiety can be considered a relatively stable characteristic of the
personality, a behavioural attitude, which reflects the manner in which the
subject tends to perceive environmental stimuli and situations as dangerous or
threatening.

Subjects with high trait anxiety show a more marked reactivity to a greater
number of stimuli and are characterized, according to Cattell and Scheier, by
high arousal, ego weakness, and a tendency towards sensitivity and guilt.

These characteristics seem to indicate a sort of predisposition to anxiety, in
the sense that subjects with high trait anxiety have a greater probability, with
respect to others, of presenting state anxiety in circumstances with low
anxiety-provoking potential and/or, for equal stimuli, of experiencing higher
levels of state anxiety. For subjects with high trait anxiety, the anxious one is
the habitual mode of response to stimuli and to environmental situations, whereas
for the others it is an exceptional mode.

\subsubsection{State Anxiety}

State anxiety manifests itself as an interruption of the emotional continuum,
that is, it provokes a rupture in the emotional balance of the person; it is
expressed through a subjective sensation of tension, worry, restlessness,
nervousness and reactivity. It is associated with an activation of the autonomic
nervous system, which provokes a series of physiological activations.

High levels of state anxiety are particularly unpleasant, disturbing and even
painful, to the point of inducing the person to put into action behavioural
adaptation mechanisms aimed at putting an end to these sensations. However, these
mechanisms may fail to achieve their purpose, leaving room for other behaviours,
this time of a maladaptive type, which lead to the opposite effect, namely a
further increase of anxiety, thus starting a vicious circle of a pathological
nature.

State anxiety, therefore, beyond certain levels of severity and duration, becomes
a pathology that may take the form of an \emph{anxiety disorder} or of one of the
related conditions described in the remainder of this chapter.

In this work, the anxiety under study is pathological state anxiety, that is, the
state anxiety that becomes disproportionate to the events that arouse it, that
persists even when the events have been overcome, or that manifests itself in the
absence of objective events.

\subsection{Symptoms of Anxiety}

The symptoms that denote an anxiety disorder may generally involve one or more of
the four functional systems that are normally coordinated by our organism to
produce adaptive responses to dangerous situations: the cognitive system, the
affective system, the behavioural system and the physiological system.

Generally, in the presence of a threat, the organism reacts by activating the
different systems: the cognitive apparatus establishes the existence of a danger
and provides selective evaluations of the threat and of the resources available
to face it; the affective component acts as an accelerator of the reaction by
increasing the sense of urgency; the behavioural component activates or inhibits
action; and the physiological component acts as a somatic mobilizer.

\paragraph{Cognitive symptoms.}
The cognitive symptoms include three different types of symptom: sensory
symptoms, difficulties of thought, and conceptual symptoms.

\begin{itemize}[leftmargin=2em]
  \item \textbf{Sensory symptoms} are significant not so much because of their
        severity or their interference with sensory and perceptual function, but
        because, being little known and not immediately controllable, they tend
        to induce more sensitive people to think they are about to ``lose their
        mind''. Such symptoms include: a confused or clouded mind, visualization
        of objects in a blurred/distant way, visualization of the environment in
        a different/unreal way, sense of unreality, and hypervigilance.
  \item \textbf{Difficulties of thought} mainly concern difficulty in
        remembering important things, confusion, the inability to control one's
        thinking, difficulty of concentration, distractibility, blocking,
        difficulty in reasoning, and loss of objectivity and perspective. They
        may be produced by a variety of factors, among them the cognitive
        inhibitions that interfere with memory and the continuous coping with
        danger.
  \item \textbf{Conceptual symptoms} reflect the preoccupation with the sense of
        vulnerability and danger. Examples are the fear of losing control, the
        fear of not being able to cope with situations, the fear of physical
        injury or of death, the fear of mental disorders, the fear of negative
        evaluations, threatening visual images, and repetitive frightening
        ideation.
\end{itemize}

\paragraph{Affective symptoms.}
The affective symptoms are often the most conspicuous in anxiety disorders. The
qualitative experiences of anxiety may differ from one situation to another, and
even from time to time within the same situation, so the anxious person may often
feel irritable, impatient, ill at ease, nervous, tense, touchy, fearful,
frightened, terrified, alarmed, appalled, excited, agitated, and so on.

\paragraph{Behavioural symptoms.}
The behavioural symptoms of anxiety generally reflect either the hyperactivity of
the behavioural system or its inhibition. There may indeed be inhibition,
immobility of muscle tone, flight, avoidance, difficult speech, defective
coordination, agitation, collapse and hyperventilation.

\paragraph{Physiological symptoms.}
The physiological symptoms of anxiety reflect the readiness of the organism, in
its globality, to react in order to protect itself by activating the branch of
the autonomic nervous system. The physiological symptoms in the different anxiety
disorders may involve the functional systems summarized in
Table~\ref{tab:physiological}.

\begin{longtable}{@{}p{0.30\textwidth}p{0.62\textwidth}@{}}
  \caption{Physiological symptoms of anxiety, grouped by the bodily system they
  affect.}\label{tab:physiological}\\
  \toprule
  \textbf{System / activity} & \textbf{Symptoms} \\
  \midrule
  \endfirsthead
  \multicolumn{2}{@{}l}{\small\itshape Table~\ref{tab:physiological} (continued)}\\
  \toprule
  \textbf{System / activity} & \textbf{Symptoms} \\
  \midrule
  \endhead
  \midrule
  \multicolumn{2}{r}{\small\itshape continued on next page}\\
  \endfoot
  \bottomrule
  \endlastfoot
  Cardiovascular activity &
    palpitations, increased heart rate, increased blood pressure, weakness,
    fainting, decreased blood pressure, decreased heart rate \\
  Respiratory activity &
    breathing difficulties, pressure in the chest, lump in the throat, sensation
    of choking, laboured, rapid or shallow breathing, etc. \\
  Neuromuscular activity &
    increased reflexes, startle reaction, twitching eyelid, insomnia, spasm,
    tremor, stiffness, agitation, strained facial expression, restless pacing,
    swaying, generalized weakness, wobbly legs, clumsy movements \\
  Gastrointestinal activity &
    abdominal pain, loss of appetite, repulsion towards food, nausea, heartburn,
    abdominal discomfort, vomiting \\
  Urinary-tract activity &
    urge to urinate, frequency of urination \\
  Skin activity &
    flushing of the face or pallid face, localized or diffuse sweating, hot or
    cold flushes, itching \\
\end{longtable}

\subsection{The DSM-5-TR Classification}

The reference nosology adopted in this work is the fifth edition, text revision,
of the \emph{Diagnostic and Statistical Manual of Mental Disorders}
(DSM-5-TR), published by the American Psychiatric Association in
2022~\cite{dsm}. With respect to the earlier DSM-IV-TR, on which the first version
of PESAD was based, the DSM-5-TR introduced a substantial reorganization of the
conditions historically grouped together as ``anxiety disorders''. The single
DSM-IV-TR class is now split across \textbf{three} separate chapters:

\begin{itemize}[leftmargin=2em]
  \item \textbf{Anxiety Disorders};
  \item \textbf{Obsessive-Compulsive and Related Disorders}; and
  \item \textbf{Trauma- and Stressor-Related Disorders}.
\end{itemize}

The key reorganizations introduced by this split, and relevant to the present
system, may be summarized as follows.

\begin{itemize}[leftmargin=2em]
  \item Obsessive-Compulsive Disorder is no longer classified among the anxiety
        disorders; it heads the new \emph{Obsessive-Compulsive and Related
        Disorders} chapter, together with body dysmorphic disorder, hoarding
        disorder, trichotillomania and excoriation disorder.
  \item Posttraumatic Stress Disorder and Acute Stress Disorder are likewise no
        longer anxiety disorders; they belong to the new \emph{Trauma- and
        Stressor-Related Disorders} chapter, which also requires an identifiable
        traumatic or stressful event as part of its diagnostic criteria.
  \item Panic disorder and agoraphobia, formerly coded jointly (as ``panic
        disorder with/without agoraphobia'' and ``agoraphobia without history of
        panic disorder''), are now \textbf{two separate diagnoses}, each of which
        can be coded independently and both of which may be assigned together.
        The \emph{panic attack} itself is no longer a codable disorder but a
        \textbf{cross-cutting specifier} that can be appended to any diagnosis
        (e.g.\ ``PTSD with panic attacks'').
  \item Social phobia was renamed \textbf{social anxiety disorder}.
  \item \textbf{Hoarding disorder} and \textbf{excoriation (skin-picking)
        disorder} are entirely new diagnoses, not present in DSM-IV-TR.
  \item \textbf{Prolonged grief disorder} is a new diagnosis added specifically
        in the DSM-5-TR text revision of 2022; it was not present in DSM-IV-TR,
        nor in the original DSM-5 of 2013.
\end{itemize}

The following three subsections describe, chapter by chapter, the disorders
actually modelled by PESAD. For each disorder a single paragraph summarizes its
core DSM-5-TR criteria, namely the key symptoms, the polythetic thresholds (of
the form ``at least $M$ of $N$'') and the required duration. The residual
\emph{other specified} and \emph{unspecified} categories of each chapter are not
modelled as positive diagnostic goals and are therefore omitted.

\subsection{Anxiety Disorders}

\subsubsection{Separation Anxiety Disorder}

The disorder is defined by developmentally inappropriate and excessive fear or
anxiety concerning separation from those to whom the individual is attached, as
evidenced by \textbf{at least 3 of 8} symptoms: recurrent excessive distress when
anticipating or experiencing separation; persistent worry about losing
attachment figures or about harm befalling them; persistent worry that an
untoward event will cause separation; reluctance or refusal to go out for fear of
separation; fear of being alone; reluctance or refusal to sleep away from home or
without a major attachment figure nearby; repeated nightmares about separation;
and repeated complaints of physical symptoms when separation occurs or is
anticipated. The fear, anxiety or avoidance must be persistent, lasting
\textbf{at least 4 weeks in children and adolescents} (younger than 18 years) and
\textbf{typically 6 months or more in adults}, and must cause clinically
significant distress or impairment.

\subsubsection{Selective Mutism}

Selective mutism is characterized by a consistent failure to speak in specific
social situations in which there is an expectation to speak (e.g.\ at school),
despite speaking in other situations. The disturbance must interfere with
educational or occupational achievement or with social communication, last
\textbf{at least 1 month} (not limited to the first month of school), and not be
attributable to a lack of knowledge of, or comfort with, the spoken language
required. It is not better explained by a communication disorder and does not
occur exclusively during an autism spectrum or psychotic disorder.

\subsubsection{Specific Phobia}

Specific phobia involves marked fear or anxiety about a specific object or
situation, which almost always provokes immediate fear and is actively avoided or
endured with intense anxiety. The fear is out of proportion to the actual danger
and to the sociocultural context, is persistent (\textbf{typically 6 months or
more}, at all ages) and causes clinically significant distress or impairment. The
disorder is coded according to the phobic stimulus, and PESAD models the
\textbf{five} DSM-5-TR stimulus types listed in Table~\ref{tab:phobia-types}.

\begin{longtable}{@{}p{0.34\textwidth}p{0.58\textwidth}@{}}
  \caption{The five stimulus types of specific phobia.}%
  \label{tab:phobia-types}\\
  \toprule
  \textbf{Type} & \textbf{Typical stimuli} \\
  \midrule
  \endfirsthead
  \toprule
  \textbf{Type} & \textbf{Typical stimuli} \\
  \midrule
  \endhead
  \bottomrule
  \endlastfoot
  Animal & spiders, insects, dogs \\
  Natural environment & heights, storms, water \\
  Blood--injection--injury & needles, invasive medical procedures, injury \\
  Situational & aeroplanes, lifts, enclosed places \\
  Other & choking or vomiting situations; in children, loud sounds or costumed
          characters \\
\end{longtable}

\subsubsection{Social Anxiety Disorder}

Social anxiety disorder is defined by marked fear or anxiety about one or more
social situations involving possible scrutiny by others (social interactions,
being observed, or performing), driven by the fear of acting in a way, or showing
anxiety symptoms, that will be negatively evaluated. The social situations almost
always provoke fear, are avoided or endured with intense anxiety, and the fear is
out of proportion to the actual threat. The disturbance is persistent
(\textbf{typically 6 months or more}) and causes clinically significant distress
or impairment. A \emph{performance only} specifier applies when the fear is
restricted to speaking or performing in public.

\subsubsection{Panic Disorder}

Panic disorder requires recurrent \textbf{unexpected} panic attacks. A panic
attack is an abrupt surge of intense fear or discomfort that peaks within
minutes, during which \textbf{at least 4 of the 13} symptoms of
Table~\ref{tab:panic-symptoms} occur. At least one attack must have been followed
by \textbf{1 month or more} of persistent concern about further attacks or their
consequences, or of significant maladaptive change in behaviour related to the
attacks. The disturbance must not be attributable to a substance or another
medical condition, nor better explained by another mental disorder.

\begin{longtable}{@{}p{0.05\textwidth}p{0.88\textwidth}@{}}
  \caption{The thirteen symptoms of a panic attack (eleven physical and two
  cognitive); at least four are required.}\label{tab:panic-symptoms}\\
  \toprule
  \# & \textbf{Symptom} \\
  \midrule
  \endfirsthead
  \toprule
  \# & \textbf{Symptom} \\
  \midrule
  \endhead
  \bottomrule
  \endlastfoot
  1 & Palpitations, pounding heart, or accelerated heart rate \\
  2 & Sweating \\
  3 & Trembling or shaking \\
  4 & Shortness of breath or smothering \\
  5 & Feelings of choking \\
  6 & Chest pain or discomfort \\
  7 & Nausea or abdominal distress \\
  8 & Feeling dizzy, unsteady, light-headed, or faint \\
  9 & Chills or heat sensations \\
  10 & Paraesthesias (numbness or tingling) \\
  11 & Derealization or depersonalization \\
  12 & Fear of losing control or ``going crazy'' \\
  13 & Fear of dying \\
\end{longtable}

\subsubsection{Agoraphobia}

Agoraphobia is defined by marked fear or anxiety about \textbf{at least 2 of the
5} situations: using public transportation; being in open spaces; being in
enclosed places; standing in line or being in a crowd; and being outside the home
alone. The situations are feared or avoided because of thoughts that escape might
be difficult, or help unavailable, should panic-like or other incapacitating or
embarrassing symptoms develop. They almost always provoke fear and are actively
avoided, require a companion, or are endured with intense anxiety. The fear is out
of proportion to the actual danger, is persistent (\textbf{typically 6 months or
more}) and causes clinically significant distress or impairment. Agoraphobia is
diagnosed irrespective of panic disorder; if both sets of criteria are met, both
diagnoses are assigned.

\subsubsection{Generalized Anxiety Disorder}

Generalized anxiety disorder is characterized by excessive anxiety and worry,
occurring \textbf{more days than not for at least 6 months}, about a number of
events or activities, which the individual finds difficult to control. The
anxiety and worry must be associated with \textbf{at least 3 of the 6} symptoms
(only \textbf{1} required in children): restlessness or feeling keyed up;
being easily fatigued; difficulty concentrating or mind going blank;
irritability; muscle tension; and sleep disturbance. The disturbance causes
clinically significant distress or impairment and is not attributable to a
substance or another medical condition.

\subsubsection{Substance/Medication-Induced Anxiety Disorder}

In this disorder, panic attacks or anxiety predominate in the clinical picture,
and there is evidence that the symptoms developed during or soon after substance
intoxication or withdrawal, or after exposure to or withdrawal from a medication
capable of producing such symptoms. The disturbance must not be better explained
by an independent anxiety disorder; persistence of symptoms for a substantial
period (\textbf{about 1 month}) after the cessation of acute withdrawal or severe
intoxication suggests an independent disorder. The symptoms must not occur
exclusively during a delirium and must cause clinically significant distress or
impairment. The condition is coded with onset during intoxication, during
withdrawal, or after medication use.

\subsubsection{Anxiety Disorder Due to Another Medical Condition}

Here panic attacks or anxiety predominate in the clinical picture, and there is
evidence from the history, physical examination or laboratory findings that the
disturbance is the direct pathophysiological consequence of another medical
condition (e.g.\ hyperthyroidism, hyperparathyroidism, phaeochromocytoma,
vestibular dysfunction, seizure disorders, or cardiopulmonary conditions). The
disturbance must not be better explained by another mental disorder, must not
occur exclusively during a delirium, and must cause clinically significant
distress or impairment.

\subsection{Obsessive-Compulsive and Related Disorders}

\subsubsection{Obsessive-Compulsive Disorder}

Obsessive-compulsive disorder is defined by the presence of obsessions,
compulsions, or both. \emph{Obsessions} are recurrent and persistent thoughts,
urges or images that are experienced as intrusive and unwanted and that cause
marked anxiety or distress, which the individual attempts to ignore, suppress or
neutralize with another thought or action. \emph{Compulsions} are repetitive
behaviours or mental acts that the individual feels driven to perform in response
to an obsession or according to rigidly applied rules, aimed at preventing or
reducing distress but not connected in a realistic way with what they are
designed to neutralize, or clearly excessive. The obsessions or compulsions must
be time-consuming (\textbf{taking more than 1 hour per day}) or cause clinically
significant distress or impairment. An \textbf{insight} specifier with three
levels (good or fair, poor, absent or delusional) and a \emph{tic-related}
specifier are available.

\subsubsection{Body Dysmorphic Disorder}

Body dysmorphic disorder consists of a preoccupation with one or more perceived
defects or flaws in physical appearance that are not observable or appear slight
to others. At some point the individual must have performed repetitive behaviours
(e.g.\ mirror checking, excessive grooming, skin picking, reassurance seeking) or
mental acts (e.g.\ comparing one's appearance with that of others) in response to
the appearance concerns. The preoccupation must cause clinically significant
distress or impairment and must not be better explained by concerns with body fat
or weight in an individual whose symptoms meet criteria for an eating disorder. A
\emph{with muscle dysmorphia} specifier and the same three-level insight specifier
apply.

\subsubsection{Hoarding Disorder}

Hoarding disorder is characterized by a persistent difficulty discarding or
parting with possessions, regardless of their actual value, owing to a perceived
need to save the items and to the distress associated with discarding them. This
difficulty results in the accumulation of possessions that congest and clutter
active living areas and substantially compromise their intended use. The hoarding
must cause clinically significant distress or impairment (including the
maintenance of a safe environment) and must not be attributable to another
medical condition or better explained by another mental disorder. A \emph{with
excessive acquisition} specifier and the three-level insight specifier apply.

\subsubsection{Trichotillomania (Hair-Pulling Disorder)}

Trichotillomania is defined by recurrent pulling out of one's hair resulting in
hair loss, together with repeated attempts to decrease or stop the pulling. The
hair pulling must cause clinically significant distress or impairment, must not
be attributable to another medical condition (e.g.\ a dermatological condition),
and must not be better explained by another mental disorder (e.g.\ attempts to
improve a perceived appearance defect in body dysmorphic disorder).

\subsubsection{Excoriation (Skin-Picking) Disorder}

Excoriation disorder consists of recurrent skin picking resulting in skin
lesions, together with repeated attempts to decrease or stop the picking. The
skin picking must cause clinically significant distress or impairment, must not
be attributable to the physiological effects of a substance (e.g.\ cocaine) or
another medical condition (e.g.\ scabies), and must not be better explained by
another mental disorder.

\subsubsection{Substance/Medication-Induced Obsessive-Compulsive and Related
Disorder}

In this disorder, obsessions, compulsions, skin picking, hair pulling, other
body-focused repetitive behaviours or other symptoms characteristic of the
chapter predominate in the clinical picture, and there is evidence that the
symptoms developed during or soon after substance intoxication or withdrawal, or
after exposure to or withdrawal from a medication capable of producing such
symptoms. The disturbance must not be better explained by an independent
obsessive-compulsive and related disorder, must not occur exclusively during a
delirium, and must cause clinically significant distress or impairment. It is
coded with onset during intoxication, during withdrawal, or after medication use.

\subsubsection{Obsessive-Compulsive and Related Disorder Due to Another Medical
Condition}

Here obsessions, compulsions, appearance preoccupations, hoarding, hair pulling,
skin picking or other characteristic symptoms predominate in the clinical
picture, and there is evidence that the disturbance is the direct
pathophysiological consequence of another medical condition. The disturbance must
not be better explained by another mental disorder, must not occur exclusively
during a delirium, and must cause clinically significant distress or impairment.
Specifiers identify the predominant symptom presentation (obsessive-compulsive,
appearance preoccupations, hoarding, hair-pulling, or skin-picking).

\subsection{Trauma- and Stressor-Related Disorders}

\subsubsection{Reactive Attachment Disorder}

Reactive attachment disorder is defined by a consistent pattern of inhibited,
emotionally withdrawn behaviour towards adult caregivers, manifested by
\textbf{both} of: the child rarely or minimally seeking comfort when distressed,
and rarely or minimally responding to comfort when distressed. It is accompanied
by a persistent social and emotional disturbance shown by \textbf{at least 2 of
3} symptoms (minimal social and emotional responsiveness; limited positive
affect; episodes of unexplained irritability, sadness or fearfulness). The child
must have experienced a pattern of extremes of insufficient care
(\textbf{at least 1 of 3}: social neglect, repeated changes of caregivers, or
rearing in settings limiting selective attachments), presumed responsible for the
disturbance. The criteria for autism spectrum disorder must not be met, the
disturbance must be evident \textbf{before age 5}, and the child must have a
developmental age of \textbf{at least 9 months}.

\subsubsection{Disinhibited Social Engagement Disorder}

This disorder is defined by a pattern of behaviour in which the child actively
approaches and interacts with unfamiliar adults, shown by \textbf{at least 2 of
4} symptoms: reduced or absent reticence in approaching unfamiliar adults; overly
familiar verbal or physical behaviour; diminished or absent checking back with a
caregiver after venturing away; and willingness to go off with an unfamiliar
adult with minimal hesitation. The behaviours must extend beyond mere
impulsivity, must follow a pattern of extremes of insufficient care
(\textbf{at least 1 of 3}, as for reactive attachment disorder) presumed
responsible for them, and the child must have a developmental age of
\textbf{at least 9 months}.

\subsubsection{Posttraumatic Stress Disorder}

Posttraumatic stress disorder requires exposure to actual or threatened death,
serious injury or sexual violence through one or more of four routes (directly
experiencing the event; witnessing it in person; learning that it happened to a
close family member or friend; or experiencing repeated or extreme exposure to
aversive details). The symptoms are organized into the four clusters of
Table~\ref{tab:ptsd-clusters}; the disturbance must last \textbf{more than 1
month} and cause clinically significant distress or impairment. A \emph{with
dissociative symptoms} specifier (depersonalization or derealization) and a
\emph{with delayed expression} specifier (full criteria not met until at least 6
months after the event) are available.

\begin{longtable}{@{}p{0.20\textwidth}p{0.10\textwidth}p{0.60\textwidth}@{}}
  \caption{The four symptom clusters of PTSD (individuals older than 6 years) and
  their polythetic thresholds.}\label{tab:ptsd-clusters}\\
  \toprule
  \textbf{Cluster} & \textbf{Threshold} & \textbf{Content} \\
  \midrule
  \endfirsthead
  \toprule
  \textbf{Cluster} & \textbf{Threshold} & \textbf{Content} \\
  \midrule
  \endhead
  \bottomrule
  \endlastfoot
  B -- Intrusion & at least 1 of 5 &
    intrusive memories; distressing dreams; dissociative reactions
    (flashbacks); psychological distress at, and physiological reactivity to,
    cues resembling the event \\
  C -- Avoidance & at least 1 of 2 &
    avoidance of distressing memories, thoughts or feelings; avoidance of
    external reminders of the event \\
  D -- Negative cognitions and mood & at least 2 of 7 &
    amnesia for the event; exaggerated negative beliefs; distorted blame;
    persistent negative emotional state; diminished interest; detachment from
    others; inability to feel positive emotions \\
  E -- Arousal and reactivity & at least 2 of 6 &
    irritable or aggressive behaviour; reckless or self-destructive behaviour;
    hypervigilance; exaggerated startle; concentration problems; sleep
    disturbance \\
\end{longtable}

For \textbf{children 6 years and younger} (the preschool subtype) the thresholds
are developmentally lowered: exposure occurs through 3 routes rather than 4, the
intrusion cluster keeps the threshold of at least 1 of 5, and avoidance and
negative cognitions are merged into a single cluster requiring \textbf{at least 1
of 6} symptoms, while the arousal cluster requires \textbf{at least 2 of 5}. The
duration (more than 1 month) and the available specifiers are unchanged.

\subsubsection{Acute Stress Disorder}

Acute stress disorder shares the same exposure criterion as PTSD (one or more of
four routes). It requires the presence of \textbf{at least 9 of 14} symptoms
drawn from five categories (intrusion, negative mood, dissociation, avoidance and
arousal), beginning or worsening after the event. The duration of the disturbance
is \textbf{3 days to 1 month after the trauma}; symptoms resolving in fewer than 3
days do not qualify, and symptoms persisting beyond 1 month lead instead to a
diagnosis of PTSD. The disturbance must cause clinically significant distress or
impairment, must not be attributable to a substance or another medical condition,
and must not be better explained by a brief psychotic disorder.

\subsubsection{Adjustment Disorder}

Adjustment disorder is the development of emotional or behavioural symptoms in
response to an identifiable stressor, occurring \textbf{within 3 months} of the
onset of the stressor. The symptoms must be clinically significant, shown by
marked distress out of proportion to the stressor and/or significant impairment.
The disturbance must not meet the criteria for another mental disorder, must not
be merely an exacerbation of a pre-existing disorder, and must not represent
normal bereavement or be better explained by prolonged grief disorder. Once the
stressor or its consequences have terminated, the symptoms do not persist for
more than a further \textbf{6 months}. The subtype is specified as: with depressed
mood; with anxiety; with mixed anxiety and depressed mood; with disturbance of
conduct; with mixed disturbance of emotions and conduct; or unspecified. Course
specifiers distinguish \emph{acute} (less than 6 months) from \emph{persistent}
(6 months or longer).

\subsubsection{Prolonged Grief Disorder}

Prolonged grief disorder, new in DSM-5-TR, requires the death of a close person
\textbf{at least 12 months ago} (\textbf{at least 6 months} in children and
adolescents). Since the death there must be a persistent grief response shown by
\textbf{one or both} of intense yearning or longing for the deceased and
preoccupation with thoughts or memories of the deceased, present nearly every day
for at least the last month. In addition, \textbf{at least 3 of 8} symptoms must
be present nearly every day for at least the last month: identity disruption; a
marked sense of disbelief about the death; avoidance of reminders that the person
is dead; intense emotional pain related to the death; difficulty reintegrating
into relationships and activities; emotional numbness; a feeling that life is
meaningless; and intense loneliness. The disturbance must cause clinically
significant distress or impairment, must clearly exceed expected social, cultural
or religious norms, and must not be better explained by another mental disorder
nor attributable to a substance or another medical condition.

\section{Identification}

In this section we describe the phases that led to the acquisition of the domain
knowledge, the identification of the requirements to be implemented, and the
definition of the reasoning strategy.


\subsection{Knowledge Acquisition}

The acquisition of the knowledge of the application domain was a complex and
laborious task. It was carried out by drawing on three different kinds of source:

\begin{itemize}
    \item \emph{Books and manuals}:
    \begin{itemize}
        \item \emph{DSM-5-TR. Diagnostic and Statistical Manual of Mental
        Disorders}~\cite{dsm}.
        \item \emph{The ABC of Psychopathology. Exploration, identification and
        treatment of mental disorders}~\cite{abc}.
    \end{itemize}

    \item \emph{Web resources}:
    \begin{itemize}
        \item \url{http://www.ipsico.org/default.htm}
        \item \url{http://www.cpsico.com/index.htm}
        \item \url{http://www.ansia-depressione.net}
    \end{itemize}

    \item \emph{Human expert}: the consultancy provided by a psychologist.
\end{itemize}

From the study of the diagnostic criteria we learned that each ``parent'' anxiety
disorder can be specialised into several ``child'' anxiety disorders.

Although the DSM-5-TR is a sound diagnostic manual, it does not allow a
knowledge engineer to encode exactly the various disorders that fall
under each category. Indeed, while it reports the diagnostic criteria related to
each ``parent'' disorder together with the possible specialisations, it does not
contain all the information needed to differentiate the diagnoses among the
``child'' disorders.

To address this issue, in addition to the use of manuals and psychology portals,
we relied on a human expert who---thanks to their experience in terms of
knowledge, heuristics and reasoning mechanisms---made it possible to simplify the
complexity of the work.

Starting from the diagnostic criteria of each ``parent'' anxiety disorder, we
carried out a process of refinement and specialisation of the various criteria so
as to identify every single ``child'' anxiety disorder. At the end of the
knowledge-acquisition work, 43 ``child'' anxiety disorders were identified; they
are reported in Table~\ref{tab:disorders}.

\begin{longtable}{p{0.34\linewidth} p{0.58\linewidth}}
\caption{Classification of the anxiety disorders handled by the expert
system.}\label{tab:disorders}\\
\toprule
\textbf{Category} & \textbf{Disorders} \\
\midrule
\endfirsthead
\toprule
\textbf{Category} & \textbf{Disorders} \\
\midrule
\endhead
\midrule
\multicolumn{2}{r}{\footnotesize\itshape continued on next page}\\
\endfoot
\bottomrule
\endlastfoot

Panic Disorder &
Panic Disorder Without Agoraphobia; Panic Disorder With Agoraphobia. \\
\addlinespace

Agoraphobia Without History of Panic Disorder &
Agoraphobia Without History of Panic Disorder. \\
\addlinespace

Specific Phobia &
Specific Phobia, ``Animal'' Type; Specific Phobia, ``Natural Environment'' Type;
Specific Phobia, ``Blood-Injection-Injury'' Type; Specific Phobia,
``Situational'' Type; Specific Phobia, ``Other'' Type. \\
\addlinespace

Social Phobia &
Generalised Social Phobia; Circumscribed Social Phobia. \\
\addlinespace

Obsessive-Compulsive Disorder &
Obsessive-Compulsive Disorder with Aggressive Content (and with Poor Insight);
with Contamination Content (and with Poor Insight); with Sexual Content (and with
Poor Insight); Hoarding Type (and with Poor Insight); with Religious Content (and
with Poor Insight); Order and Symmetry Type (and with Poor Insight); Somatic Type
(and with Poor Insight); Superstitious Type (and with Poor Insight); Perfection
and Responsibility Type (and with Poor Insight); with Miscellaneous Content (and
with Poor Insight). \\
\addlinespace

Post-Traumatic Stress Disorder &
Acute Post-Traumatic Stress Disorder; Acute Post-Traumatic Stress Disorder with
Delayed Onset; Chronic Post-Traumatic Stress Disorder; Chronic Post-Traumatic
Stress Disorder with Delayed Onset. \\
\addlinespace

Acute Stress Disorder &
Acute Stress Disorder. \\
\addlinespace

Generalised Anxiety Disorder &
Generalised Anxiety Disorder. \\
\addlinespace

Anxiety Disorder Due to a General Medical Condition &
With Generalised Anxiety; With Panic Attacks; With Obsessive-Compulsive
Symptoms. \\
\addlinespace

Substance-Induced Anxiety Disorder &
With Generalised Anxiety; With Panic Attacks; With Obsessive-Compulsive Symptoms;
With Phobic Symptoms. \\
\end{longtable}


\subsection{Requirements}

In the following sections we provide a brief description of the requirements that
the system must satisfy and of the features to be implemented.


\subsubsection{Programming Language}

The expert system must be implemented in Prolog.

Prolog, an acronym for \emph{PROgramming in LOGic}, is a very powerful and
flexible logic programming language. It is particularly well suited to Artificial
Intelligence and to non-numerical programming.

By using Prolog as the programming language, it is possible to express the
knowledge about the identified objects and relations in the form of facts and
rules. It is also possible to ask questions of the user so as to derive further
information useful to the inference process.

Two behavioural features required by the system are not provided natively by
Prolog: the handling of uncertainty and an explanation module. They are discussed
below.


\subsubsection{Handling of Uncertainty}

The first behavioural feature we want to implement is the handling of
uncertainty.

Prolog reasons in a categorical way, in the sense that every element can only be
either true or false, with no other possibility.

In the medical field, this kind of reasoning turns out to be too restrictive. The
diagnosis of disorders or diseases may, in fact, rest on rather uncertain
knowledge, both because symptoms and causes may be only probably correlated with
a disorder or disease, and because the patient may not answer with complete
certainty.

The handling of the uncertainty intrinsic to the available knowledge becomes even
more necessary in a critical domain such as the diagnosis of anxiety disorders,
where the mental disorder must be managed as quickly and correctly as possible, in
order to prevent the exacerbation of the symptoms from ``destroying'' the overall
functioning of the patient or from significantly interfering with their
occupational or social functioning.

From these premises, the need emerges to move away from the categorical nature of
assertions and to converge towards a probabilistic nature, in which it is possible
to assign a degree of certainty to every known fact or rule.

This notion of certainty is not to be understood in the mathematical sense of the
term, but rather as the subjective confidence one has in the truth of an assertion
or in the fact that an inference rule turns out to be applicable.

Indeed, it is still a matter of debate whether it is more appropriate to use
probability theory in its most rigorous sense or in a more informal and
simplified version.

In the project under examination, we chose to use the more informal and
simplified version for two simple reasons. First, because we want the expert
system to emulate human reasoning, which is certainly intuitive and not very
mathematical. Second, because one cannot disregard the greater computational
simplicity of an intuitive and less rigorous method.


\subsubsection{Explanation Module}

The second behavioural feature concerns the implementation of an explanation
module.

In an expert system, the explanation capability is of fundamental importance.
Users, in fact, tend not to trust an artificial system slavishly, nor to accept
blindly the solution of a problem produced in an uncontrollable way.

This capability is even more necessary in uncertain domains, such as that of
medical diagnosis, in order to increase the trust that the physician or the expert
places in the advice of the system, and to enable them to discover a possible
error in the reasoning process.

In keeping with this requirement, the expert system under examination provides
explanations about its own way of proceeding on two occasions: when it is asked
\emph{why} a given question is being posed, and when it is asked \emph{how} the
solution of a problem was reached.


\subsection{Reasoning Strategy}

As far as the type of inference to be used is concerned, the choice fell on
backward chaining.

Backward chaining is a type of inference in which one starts from a possible
solution of the problem and tries to verify, by applying the production rules in
reverse, whether there exists a rule capable of providing that solution. It is
opposed to forward chaining, in which one starts from the known facts contained in
the knowledge base and applies the production rules forward in order to deduce new
facts. Like the handling of uncertainty, backward chaining is not provided
natively by Prolog and must be implemented by the inference engine.

The inference process will moreover use the \emph{left-most} computation rule and
the \emph{depth-first} search strategy. Thus, in the inference process, given a
query

\[
\texttt{:-}\ G_1,\, G_2,\, \dots,\, G_n
\]

the left-most literal will always be selected. Given the literal to be resolved
(for example, $G_1$), the first clause whose head is unifiable with $G_1$ will
always be selected.

As for the behaviour that the inference process must adopt upon the occurrence of
possible negative answers, we opted for a ``non-failure'' policy. Considering, in
fact, that the expert system has to handle the uncertainty of the information, we
decided to process negative answers as positive answers, but with a
``complemented'' degree of certainty. For instance, a negative answer with a
degree of certainty of $0.7$ for a feature required as true will be processed by
the system as a positive answer with a degree of certainty of $0.3$. In this way,
it is possible to bring every computation to completion, returning the degree of
certainty with which the conclusion turns out to be true.

Finally, since some diagnostic criteria require the verification of undesirable
features, the system must provide the possibility of specifying, for a given goal,
which sub-goals are required as true and which as false. Depending on the case, it
will be the inference engine that correctly computes the degree of certainty.

The reasoning strategy introduced in this section will be taken up again and
documented in depth in the sections devoted to the description of the inference
engine.

\section{Conceptualization}

In this section we present an analysis of the process that led to expressing the
knowledge in logical form. In particular, we discuss the objects that were
identified and the relations and properties of those objects.

The knowledge is organised over four inference levels, which structure the
reasoning from the most abstract diagnosis down to the elementary observations:

\begin{enumerate}[label=(\Roman*)]
    \item \emph{Disorders} ---the ``child'' anxiety disorders that the system is
    able to diagnose, encoded by the relation \texttt{disorder/3};
    \item \emph{Sub-disorders and types} ---the ``parent'' or
    ``non-codifiable'' disorders and the specialisation types that, combined with
    them, define a disorder, encoded by the relation \texttt{disorder\_nc/4};
    \item \emph{Symptomatic manifestations and features} ---the complex symptoms
    and the features attached to them, encoded by the relations
    \texttt{symptomatic\_manifestation/4} and \texttt{symptomatic\_feature/4};
    \item \emph{Elementary symptoms} ---the atomic facts reported by the user,
    encoded by the relation \texttt{symptom/3}.
\end{enumerate}


\subsection{Objects}

The conceptualization phase brought several objects to light. These objects can be
classified into categories.

In Table~\ref{tab:objects}, rather than weighing down the discussion with a list
of every single identified object, we prefer to report the categories to which the
objects belong, focusing attention on the description of what they represent. For
each category, a few examples of objects are also given.

\begin{longtable}{p{0.22\linewidth} p{0.42\linewidth} p{0.28\linewidth}}
\caption{Objects identified during the conceptualization
phase.}\label{tab:objects}\\
\toprule
\textbf{Object Category} & \textbf{Description} & \textbf{Examples} \\
\midrule
\endfirsthead
\toprule
\textbf{Object Category} & \textbf{Description} & \textbf{Examples} \\
\midrule
\endhead
\midrule
\multicolumn{3}{r}{\footnotesize\itshape continued on next page}\\
\endfoot
\bottomrule
\endlastfoot

Disorder &
The ``child'' anxiety disorders that the expert system is able to diagnose. &
Panic Disorder Without Agoraphobia; Specific Phobia, Animal Type; Generalised
Anxiety Disorder; \dots \\
\addlinespace

Non-Codifiable Disorder &
Both the ``parent'' anxiety disorders to which the ``child'' disorders refer and
the so-called ``non-codifiable'' disorders, that is, those disorders that, without
a precise specialisation or placement, cannot be recorded. &
Agoraphobia; Specific Phobia; Post-Traumatic Stress Disorder; Panic Disorder;
Social Phobia; \dots \\
\addlinespace

Type of Non-Codifiable Disorder &
The various specialisations that, combined with the non-codifiable disorders,
allow a disorder to be defined. &
Generalised Type; Circumscribed Type; Animal Type; Natural Environment Type;
\dots \\
\addlinespace

Symptomatic Manifestation &
The various manifestations of symptoms that denote an anxiety disorder.
Typically, a symptomatic manifestation is made up of several criteria (complex
symptoms). Each criterion is composed of one or more symptoms, which need not all
be mandatorily required for the diagnosis: it may be the case that a minimum
number of symptoms out of a larger set is required. &
Increased Arousal; Psychological Trauma; Social Phobic Anxiety; Agoraphobic
Anxiety; Unexpected Panic Attack; Generalised Anxiety; \dots \\
\addlinespace

Symptomatic Feature &
The features connected to the symptoms that allow one disorder to be determined or
distinguished from another. &
Duration of at least 6 months; Frequent panic attacks; Onset of symptoms within 1
month of the traumatic event; \dots \\
\addlinespace

Elementary Symptom &
The atomic symptoms that, once assembled, give rise to the symptomatic
manifestations. &
Depersonalisation; Terror; Avoidance; Nausea; \dots \\
\addlinespace

Patient &
The categories of patient to whom the diagnosis is addressed. Placing the patient
in one of the three age ranges makes it possible to distinguish both the number of
symptomatic manifestations or features required by a given disorder and the number
of symptoms expected from a single symptomatic manifestation. &
Child (0--9 years); Adolescent (10--17 years); Adult (18+). \\
\addlinespace

Degree of Insight of the Obsessions &
The degrees of insight into the reasonableness of the obsessions-compulsions. &
High Insight; Low Insight; \dots \\
\addlinespace

Type of Anxiety &
The types of anxiety. The distinction is based on the main reason that triggers
the anxious process. &
Anxiety due to embarrassment in public; Anxiety due to concern for one's own
health; \dots \\
\addlinespace

Type of Panic Attack &
The two types into which a panic attack can be specialised. &
Unexpected Panic Attack; Situationally Triggered Panic Attack. \\
\addlinespace

Type of Phobic Reaction &
The types of reaction that an individual may have towards a feared situation or
object. &
Endurance with distress; Avoidance; \dots \\
\addlinespace

Symptom Duration &
The various durations used to verify all the symptoms that require the presence of
the symptoms for a precise length of time. &
Less than 2 days; 2 days -- less than 1 month; 1 month -- less than 3 months;
3 months -- less than 6 months; 6 months or more. \\
\addlinespace

Symptom Onset &
The various periods within which the acute stress disorder and the post-traumatic
stress disorder can arise. &
Less than 1 month; 1 month -- less than 6 months; 6 months or more. \\
\addlinespace

Panic Attack Frequency &
The various frequencies of panic attacks required in panic disorder. &
1; 2--5; 6+. \\
\end{longtable}


\subsection{Relations and Properties}

Table~\ref{tab:relations} reports the relations and the properties of the objects.
The first five relations correspond to the four inference levels described at the
beginning of this section.

\begin{longtable}{p{0.46\linewidth} p{0.46\linewidth}}
\caption{Relations and properties of the objects.}\label{tab:relations}\\
\toprule
\textbf{Relation} & \textbf{Description} \\
\midrule
\endfirsthead
\toprule
\textbf{Relation} & \textbf{Description} \\
\midrule
\endhead
\midrule
\multicolumn{2}{r}{\footnotesize\itshape continued on next page}\\
\endfoot
\bottomrule
\endlastfoot

\texttt{patient(Patient)} &
Indicates the type of patient to whom the diagnosis is addressed. \\
\addlinespace

\texttt{disorder/3} &
States that the given object is a ``child'' disorder and is diagnosable for the
specified patient (level~I). \\
\addlinespace

\texttt{disorder\_nc/4} &
States that the given object is a ``parent'' or ``non-codifiable'' disorder,
together with its specialisation type, and is diagnosable for the specified
patient (level~II). \\
\addlinespace

\texttt{symptomatic\_manifestation/4} &
States that the given object is a symptomatic manifestation and is diagnosable for
the specified patient (level~III). \\
\addlinespace

\texttt{symptomatic\_feature/4} &
States that the given object is a symptomatic feature and is diagnosable for the
specified patient (level~III). \\
\addlinespace

\texttt{symptom/3} &
Encodes a criterion (complex symptom) at the elementary level (level~IV): the
list argument indicates the various elementary symptoms that compose the
criterion, while the numeric argument indicates the number of symptoms mandatorily
required to diagnose the criterion. \\
\addlinespace

\texttt{aetiological\_evidence([Evidence\_List], Evidence\_Number)} &
Indicates an aetiological evidence, where \texttt{Evidence\_List} lists the
various ways of verifying the aetiological link between the disorder and the
factor under examination (drug, withdrawal, intoxication), and
\texttt{Evidence\_Number} indicates the number of evidences mandatorily required
for the diagnosis. \\
\addlinespace

\texttt{phobic\_reaction\_type([Options])} &
Indicates a type of phobic reaction, where \texttt{Options} lists the admissible
reaction options. \\
\end{longtable}

\section{Implementation}

This section describes how PESAD is realised: the tools adopted, the modular
architecture of the system, the representation of the knowledge base, the
mechanisms of the inference engine ---knowledge acquisition, uncertainty
management and explanation--- and the bilingual user interface.

\subsection{Tools Used}

PESAD is implemented in \emph{SWI-Prolog}, a mature and free implementation of
the Prolog logic-programming language. Prolog was chosen because the diagnostic
knowledge is naturally expressed as \emph{rules} (definite clauses) and because
its resolution mechanism provides backward chaining and backtracking out of the
box. The features that Prolog does not provide natively ---reasoning under
uncertainty and the explanation of the reasoning process--- are implemented
explicitly on top of the language, as described below.

\subsection{System Architecture}

The system is organised by \emph{responsibility} rather than as a single
monolithic file. This separation mirrors the conceptual architecture of an
expert system (knowledge base, inference engine, user interface) and makes each
concern independently readable and maintainable. Table~\ref{tab:modules} lists
the modules and the directives that load them are collected in
\texttt{start.pl}.

\begin{table}[h]
\centering
\caption{Modular organisation of PESAD.}
\label{tab:modules}
\small
\begin{tabular}{@{}lp{9.2cm}@{}}
\toprule
\textbf{Module} & \textbf{Responsibility} \\
\midrule
\texttt{start.pl} & Bootstrap: directives and loading of the modules. \\
\texttt{shell.pl} & User shell: REPL loop, language selection, command dispatch, presentation of results, and the \texttt{decode/2}\,/\,\texttt{explain/2} localization dispatch. \\
\texttt{ask.pl} & Knowledge acquisition: the \emph{criterion}, \emph{multiple} and \emph{selective} question protocols, with answer parsing and certainty elicitation. \\
\texttt{uncertainty.pl} & Certainty-factor calculus, certainty-based ordering and inference-level index bookkeeping. \\
\texttt{explanation.pl} & Printing of the proof tree (\emph{how} a conclusion was reached). \\
\texttt{memory.pl} & Working-memory management (facts and proof trees). \\
\texttt{utils.pl} & General-purpose helpers (output spacing, list operations). \\
\texttt{disorders.pl} & Knowledge base: codable Axis~I disorders (diagnostic goals) and entry points. \\
\texttt{criteria.pl} & Knowledge base: sub-disorders, types, symptomatic manifestations and features, and the questions. \\
\texttt{interface\_en.pl}, \texttt{interface\_it.pl} & English / Italian localization (\texttt{label/3}, \texttt{explanation/3}, \texttt{command/3}). \\
\bottomrule
\end{tabular}
\end{table}

\subsubsection{Knowledge Base}

The knowledge base is structured over \emph{four inference levels} that descend
from the abstract diagnoses to the concrete observations:

\begin{enumerate}[label=(\Roman*)]
    \item \textbf{Disorders} ---the codable Axis~I diagnoses, i.e. the goals of
    the system (\texttt{disorder/3});
    \item \textbf{Sub-disorders and types} ---non-codable building-block
    disorders and their specifiers (\texttt{disorder\_nc/4} and the
    \texttt{*\_type/4} predicates);
    \item \textbf{Symptomatic manifestations and features}
    (\texttt{symptomatic\_manifestation/4}, \texttt{symptomatic\_feature/4});
    \item \textbf{Elementary symptoms / facts} elicited from the user
    (\texttt{symptom/3}).
\end{enumerate}

Each rule carries, besides the patient and the resulting certainty factor, the
truth value it is asserted with. As an example, the rule for \emph{Panic
Disorder} requires the presence of the panic-disorder criteria and the
\emph{absence} of an associated substance-induced alteration and of an
associated general medical condition (the shared exclusion):

\begin{lstlisting}
disorder(panic_disorder, Patient, CF) :-
    disorder_nc(panic_disorder, Patient, CF1, true),
    exclude_substance_and_medical(Patient, CF2, CF3),
    certainty(panic_disorder, 1, [CF1,CF2,CF3], 0.97, CF, true).
\end{lstlisting}

The last argument of \texttt{certainty/6} is the truth value, the constant
\texttt{0.97} is the certainty of the rule itself, and the integer \texttt{1}
identifies the inference level, which is used by the explanation module to build
the proof tree.

\subsubsection{Inference Engine}

The inference engine is the part of the system that, given the facts stated by
the user, derives the diagnoses. It comprises three concerns: the acquisition of
information from the user, the management of uncertainty, and the explanation of
the reasoning.

\paragraph{Knowledge Acquisition.}
Questions are posed to the user through three protocols, implemented in
\texttt{ask.pl}.

\emph{Criterion ask} verifies a criterion made of one or more items, each a
yes/no question; a criterion of $N$ items may require that only $M$ of them
(with $1 \le M \le N$) be satisfied. Its certainty is obtained by taking the $M$
items with the highest certainty and then keeping the \emph{lowest} certainty
among them:
\begin{equation}
\cf_{\text{criterion}} \;=\; \min \big( \text{top-}M \{\, \cf_1, \dots, \cf_N \,\} \big).
\end{equation}
For instance, with $N=3$ items of certainties $\{0.4, 0.7, 0.6\}$ and $M=2$, the
two highest are $\{0.7, 0.6\}$ and the resulting certainty is $0.6$. If the user
answers \emph{no} to an item with certainty $c$, the item contributes the
complementary certainty $1-c$.

\emph{Multiple ask} lets the user choose one value from a list of options. If
the chosen answer is among those admitted by the question, the certainty equals
the one declared by the user; otherwise the residual certainty is distributed
over the remaining options:
\begin{equation}
\cf \;=\;
\begin{cases}
\cf_{\text{answer}} & \text{if the answer is admitted,}\\[4pt]
\dfrac{1-\cf_{\text{answer}}}{\,|\text{options}|-1\,} & \text{otherwise.}
\end{cases}
\end{equation}

\emph{Selective ask} is used for single, mutually exclusive choices (e.g. the
patient's age group): the admitted option succeeds with certainty $1.0$, while
any other choice makes the goal fail.

\paragraph{Uncertainty Management.}
Reasoning under uncertainty is implemented with \emph{certainty factors}
(\cf{}~$\in [0,1]$), managed in \texttt{uncertainty.pl}. The certainty of the
\emph{antecedent} of a rule, whose premises must all hold (a conjunction), is
the minimum of the premises' certainties:
\begin{equation}
\cf_{\text{ant}} \;=\; \min_i \cf_i .
\end{equation}
The certainty of the \emph{consequent} combines the antecedent certainty with
the certainty of the rule. For a goal asserted as true the two are multiplied;
for a negated goal the complement of the product is taken:
\begin{equation}
\cf_{\text{cons}} \;=\;
\begin{cases}
\cf_{\text{rule}} \cdot \cf_{\text{ant}} & \text{(positive goal),}\\[4pt]
1 - \cf_{\text{rule}} \cdot \cf_{\text{ant}} & \text{(negated goal).}
\end{cases}
\end{equation}
These two predicates are at the heart of the calculus:

\begin{lstlisting}
antecedent_certainty(Antecedent_CF_List, Antecedent_CF) :-
    min_list(Antecedent_CF_List, Antecedent_CF).

consequent_certainty(Goal, 1, Antecedent_CF, Rule_CF, Consequent_CF, true) :-
    Consequent_CF is Rule_CF * Antecedent_CF,
    ...
consequent_certainty(Goal, 1, Antecedent_CF, Rule_CF, Consequent_CF, false) :-
    Consequent_CF is 1 - (Rule_CF * Antecedent_CF),
    ...
\end{lstlisting}

While computing the consequent, each rule is asserted (as \texttt{rule\_I/4},
\texttt{rule\_II/5} or \texttt{rule\_III/6}, according to its inference level)
together with a hierarchical index. These assertions form the proof tree that
the explanation module later traverses.

\paragraph{Certainty factors versus fuzzy logic.}
The conjunction of the premises is, in fact, a \emph{fuzzy} operation: PESAD
evaluates it with a configurable t-norm, selected at start-up and stored in the
dynamic fact \texttt{uncertainty\_method/1}. Two methods are provided:
\begin{itemize}
  \item \texttt{cf} (the default) uses the \emph{minimum} (G\"odel t-norm),
  i.e. the certainty of a criterion equals that of its \emph{weakest} satisfied
  premise: $\cf_{\text{ant}} = \min_i \cf_i$;
  \item \texttt{fuzzy} uses the \emph{product} t-norm, i.e. the algebraic
  product of the premises: $\cf_{\text{ant}} = \prod_i \cf_i$.
\end{itemize}
The remaining steps (rule chaining $\cf_{\text{rule}}\cdot\cf_{\text{ant}}$ and
the negation $1-x$) are shared. The two methods are compared on a suite of
clinical vignettes in Chapter~\ref{chap:validation}; both rank the correct
diagnosis first, but the product t-norm systematically \emph{under-rates
polythetic disorders}: a fully met post-traumatic stress disorder, which
combines seven symptom clusters,
collapses to $28\%$ because every premise multiplies the result, whereas the
minimum keeps the confidence at a clinically plausible $66\%$ regardless of how
many criteria a disorder has. For a diagnostic task in which a syndrome is ``as
strong as its weakest satisfied criterion'', the minimum is the more faithful
operator, and it is therefore kept as the default; the product t-norm remains
available for comparison and experimentation.

A richer fuzzy treatment is also provided as a third method
(\texttt{fuzzy\_membership}, in \texttt{fuzzy.pl}): for the graded clinical
variables ---symptom duration, panic-attack frequency and degree of insight---
the clinician enters a \emph{crisp} value (e.g. the duration in months) instead
of choosing a discrete option, and a trapezoidal membership grade becomes the
certainty. This replaces the hard diagnostic cut-offs with smooth boundaries:
for instance, against the ``at least six months'' criterion a five-month
duration scores $0.67$ rather than abruptly failing, while the conjunction is
still aggregated with the minimum. Table~\ref{tab:duration-membership}
illustrates the membership of the ``at least six months'' fuzzy set.

\begin{table}[h]
\centering
\caption{Membership grade of the ``duration at least 6 months'' fuzzy set.}
\label{tab:duration-membership}
\small
\begin{tabular}{@{}lcccccc@{}}
\toprule
\textbf{duration (months)} & 2 & 3 & 4 & 5 & 6 & 8 \\
\midrule
\textbf{membership} & 0\% & 0\% & 33\% & 67\% & 100\% & 100\% \\
\bottomrule
\end{tabular}
\end{table}

\paragraph{Explanation Module.}
A good expert system must justify its behaviour. PESAD does so in two ways,
implemented in \texttt{ask.pl} (the \emph{why}) and \texttt{explanation.pl} (the
\emph{how}).

\emph{Why} questions: while answering a yes/no or multiple-choice question, the
user may reply \texttt{why} (\texttt{perche}) to receive an explanation of the
diagnostic reason behind the question; the same question is then asked again.

\emph{How} questions: after a diagnosis the user may ask how it was reached. The
predicate \texttt{print\_tree/1} walks the asserted rules across the three
inference levels and prints, indented by level, each (sub-)conclusion with its
certainty, distinguishing what was \emph{deduced} (intermediate rules) from what
was \emph{stated} by the user (leaf facts).

\subsubsection{User Interface}

The user interacts with the system through a textual shell (\texttt{shell.pl})
that implements a read--evaluate--print loop. At start-up the user selects the
interface language; the dynamic fact \texttt{language/1} then drives every
message through the localization dispatch:

\begin{lstlisting}
decode(Key, Text) :- language(Language), label(Language, Key, Text), !.
decode(Key, Key).
\end{lstlisting}

All user-visible text is therefore stored as language-neutral \emph{keys} that
each locale file (\texttt{interface\_en.pl}, \texttt{interface\_it.pl}) maps to a
\texttt{label/3} string and, where relevant, to an \texttt{explanation/3} string.
The command vocabulary is localized in the same way through \texttt{command/3},
so that the user may type either the English or the Italian command words while
the engine works on a single set of internal command atoms
(\texttt{investigation}, \texttt{control}, \texttt{facts}, \texttt{help},
\texttt{clean}, \texttt{quit}). Adding a new language amounts to providing a new
locale file, with no change to the engine or the knowledge base.

\chapter{Use and Functionality of the System}

This chapter provides a brief description of how the system is started and of the
functionalities it implements.


\section{Starting the System}

In order to run the expert system it is sufficient to copy the \texttt{PESAD}
folder, which contains the modules that implement the system, into a directory of
choice. The system relies on the SWI-Prolog environment and is launched by
consulting the start-up file \texttt{start.pl}, located in the \texttt{PESAD}
folder, directly from the command line:

\begin{lstlisting}
$ swipl start.pl
\end{lstlisting}

The very first action performed at start-up is the selection of the interface
language. The system asks the user to choose between English and Italian, and the
whole session is then conducted in the selected language:

\begin{lstlisting}
Select language / Seleziona lingua:
	1: English
	2: Italiano
 > 1.
\end{lstlisting}

Once the language has been chosen, the system prints a short presentation message
followed by the list of available commands, and then waits for the user to type a
command:

\begin{lstlisting}
-------------------------------------------------------------------
-- PESAD - Prolog Expert System for Anxiety,
--         Obsessive-Compulsive and Trauma-Related Disorders
--
--	knowledge base: DSM-5-TR
--	author:         Donato Meoli
-------------------------------------------------------------------

Available commands:

	investigation.	- Run the full investigation of all disorders;
	control.	- Run the check of a single disorder;
	facts.		- Show the asserted facts;
	help.		- Show the list of available commands;
	clean.		- Clear the memory;
	quit.		- Close the interface.

>
\end{lstlisting}

At this point the system remains idle, waiting for a user command. Once a command
has been typed, the system performs a short validity check on it: if the command
is not recognised, an error message is shown and the system waits again for a
valid command. As a general rule, every input typed by the user is followed by
rigorous validity checks.


\section{System Commands}

The expert system provides six different commands:

\begin{itemize}
  \item \texttt{investigation} --- starts the complete investigation of all the
        disorders represented in the system;
  \item \texttt{control} --- starts the check of a single disorder, chosen among
        those implemented in the system;
  \item \texttt{facts} --- displays the facts temporarily asserted by the user
        during the current consultation;
  \item \texttt{help} --- displays the list of commands accepted by the system;
  \item \texttt{clean} --- clears the working memory;
  \item \texttt{quit} --- closes the user interface.
\end{itemize}

The command words are always in English, regardless of the selected interface
language: only the questions, diagnoses and explanations that the user reads are
localised, whereas the commands are control tokens and are kept uniform.

Throughout a consultation, when the inference engine lacks the information needed
to continue its reasoning, it poses questions to the user. Symptom questions are
answered with \texttt{yes} or \texttt{no}; in addition, the user may type
\texttt{why} (or \texttt{perche} in the Italian interface) to obtain an
explanation of why the question is being asked. After a \texttt{yes}/\texttt{no}
answer, the system asks for a certainty value, which must be a real number in the
interval $[0.0, 1.0]$ expressing how confident the user is in the given answer.


\subsection{Investigation}

The \texttt{investigation} command launches the procedure that investigates the
symptoms affecting a given patient, in order to diagnose a list of possible
anxiety disorders. The result of the investigation is therefore an \emph{ordered}
list of disorders that are candidates to be the cause of the patient's distress.

This functionality is characterised by two design choices. The first is to
structure the reasoning so that, once an Axis~I disorder has been resolved, the
system does not stop the exploration but continues to examine alternative paths.
In this way the inference engine produces a list of anxiety disorders, each
associated with a certainty score. This score, obtained as the product of the
certainty of the premise and the certainty of the rule representing the disorder,
allows a \emph{ranking} of the anxiety disorders, from the most certain to the
least certain.

The second choice is to report in the result list all the disorders that were
examined, together with their certainty degree. This was considered appropriate
because the expert user of the system can benefit from a complete view of the
anxiety disorders and, whenever a discrepancy with their own reasoning is
observed, can investigate the reason for the difference by displaying the proof
tree of the disorders involved.

In the code, this functionality is implemented through the following predicate:

\begin{lstlisting}[language=Prolog]
solve(investigation) :-
    delete_trees,
    investigate_goals(Goals_List),
    descending_order_goals(Goals_List, Sorted_Goals_List),
    display_investigation(Sorted_Goals_List),
    investigation_how(Sorted_Goals_List).
\end{lstlisting}

The predicate \texttt{delete\_trees} removes any trace of the chain of rules and
facts that led to the various goals in previous executions. This operation does
\emph{not} remove the facts dynamically asserted by the user during previous
consultations: the two operations are deliberately kept separate, so that all the
facts stated by the user remain in memory across subsequent executions, while the
traces of the rules and facts used in the proof trees are removed at every run.
This choice prevents partial traces of disorders from remaining in memory should
the investigation procedure terminate prematurely because of unexpected errors or
incorrect use of the system. If the user wishes to clear the entire working
memory, the \texttt{clean} command can be used.

The predicate \texttt{investigate\_goals/1} first determines the patient's age
class --- child, adolescent or adult --- and then, for every disorder
represented, verifies the required features, asking the user questions whenever
the available information is insufficient. The predicate
\texttt{descending\_order\_goals/2} sorts the resolved goals in descending order
of certainty, \texttt{display\_investigation/1} prints the resulting list of
disorders together with their certainty degree, and
\texttt{investigation\_how/1} manages the optional request for an explanation of
how a conclusion was reached. A single derivation tree is printed at a time, to
make it easier to read: after the diagnosis has been shown, the system asks the
user whether they want the ``how''; if the answer is \texttt{yes}, the system
asks for the code of the disorder for which the explanation should be produced,
prints the corresponding tree, and then allows the user either to leave the
investigation procedure or to request further explanations for other disorders.

A concise example of an investigation session is shown below:

\begin{lstlisting}
> investigation.

[ ... the system asks the relevant symptom questions ... ]

Does the patient experience recurrent and unexpected panic attacks?
? (yes, no) > yes.
Enter the certainty degree (from 0.0 to 1.0) > 0.9.

Does the patient avoid situations from which escape might be difficult?
? (yes, no) > why.
Pervasive fear or avoidance of situations from which escape might be
difficult would direct the diagnosis in favor of agoraphobia.
? (yes, no) > yes.
Enter the certainty degree (from 0.0 to 1.0) > 0.8.

========================== DIAGNOSIS ==========================

1) Panic Disorder with a certainty degree of 81%

2) Agoraphobia with a certainty degree of 64%

Do you want to know how they were deduced? (yes, no) > yes.

Enter the code corresponding to the chosen option > 1.

[ ... the proof tree for Panic Disorder is printed ... ]

Do you want to know how they were deduced? (yes, no) > no.
\end{lstlisting}


\subsection{Control}

The \texttt{control} command launches the procedure that checks whether the
patient suffers from a given anxiety disorder. This functionality is of
fundamental importance whenever the expert has already carried out a preliminary
analysis of the patient, or has data that allow a particular type of disorder to
be hypothesised: with this tool the expert can confirm or reject that hypothesis.

In the code, the control of a given disorder is implemented through the following
predicate:

\begin{lstlisting}[language=Prolog]
solve(control) :-
    selection_goal(Goal_Name),
    delete_tree(Goal_Name),
    solve_goal(Goal_Name, CF),
    display_control(Goal_Name, CF),
    control_how(Goal_Name).
\end{lstlisting}

The predicate \texttt{selection\_goal/1} lets the user select the disorder to be
checked from a list of possible options, retrieved from the up-to-date list of
disorders present in the knowledge base. The predicate
\texttt{delete\_tree/1} then removes any trace of the chain of rules and facts
that led to the conclusion of that goal in a previous execution; as in the
investigation case, this operation removes only the proof-tree traces and not the
facts asserted by the user, again to avoid leaving partial traces in memory if
the procedure terminates prematurely. The entire working memory can still be
cleared with the \texttt{clean} command.

The predicate \texttt{solve\_goal/2} is very similar to
\texttt{investigate\_goals/1}, the difference being that the inference engine
verifies exclusively the disorder unifying with \texttt{Goal\_Name}; as before,
the reasoning is driven by the patient's age class, and questions are posed to the
user only when new information is needed. The predicate returns the certainty
degree with which the disorder is found to be true. The predicate
\texttt{display\_control/2} prints the single result of the check, and
\texttt{control\_how/1} asks the expert whether they are interested in an
explanation of how the conclusion was reached: if the answer is \texttt{yes}, the
system immediately prints the corresponding proof tree; otherwise it closes the
control procedure and waits for new commands.

A concise example of a control session is shown below:

\begin{lstlisting}
> control.

Which of the following disorders do you want to check?

	1: Panic Disorder
	2: Agoraphobia
	3: Generalized Anxiety Disorder
	[ ... ]

Enter the code corresponding to the chosen option > 1.

[ ... the system asks only the questions relevant to Panic Disorder ... ]

Does the patient experience recurrent and unexpected panic attacks?
? (yes, no) > yes.
Enter the certainty degree (from 0.0 to 1.0) > 0.9.

========================== DIAGNOSIS ==========================

The patient suffers from Panic Disorder with a certainty degree of 81%

Do you want to know how it was deduced? (yes, no) > yes.

[ ... the proof tree for Panic Disorder is printed ... ]
\end{lstlisting}


\subsection{Facts}

The \texttt{facts} command displays the facts that the user has temporarily
asserted through the interface during the current consultation. Each fact is
printed in the form \texttt{(Answer, Attribute, Value, Certainty)}, where:

\begin{itemize}
  \item \texttt{Answer} is the type of response, \texttt{yes} or \texttt{no}. For
        facts asserted through binary (yes/no) questions, the value is
        \texttt{yes} when the patient answered affirmatively and \texttt{no} when
        the answer was negative. For facts asserted through multiple-choice
        questions the value is always \texttt{yes}, since the answer corresponds
        to the option selected by the user;
  \item \texttt{Attribute} is the symptom or feature being investigated;
  \item \texttt{Value} is the item or option associated with the answer;
  \item \texttt{Certainty} is the certainty value, in $[0.0, 1.0]$, that the user
        associated with the answer.
\end{itemize}

An example of the output produced by this command is the following:

\begin{lstlisting}
> facts.

Asserted facts:

(yes, panic_attacks, recurrent_unexpected, 0.9)
(no, medical_condition, present, 0.95)
(yes, avoidance, social_situations, 0.8)
\end{lstlisting}


\subsection{Help}

The \texttt{help} command reprints the menu listing all the commands accepted by
the system, together with a short description of each, exactly as shown at
start-up. It is useful as a quick reminder of the available functionalities
during a consultation:

\begin{lstlisting}
> help.

Available commands:

	investigation.	- Run the full investigation of all disorders;
	control.	- Run the check of a single disorder;
	facts.		- Show the asserted facts;
	help.		- Show the list of available commands;
	clean.		- Clear the memory;
	quit.		- Close the interface.
\end{lstlisting}


\subsection{Clean}

The \texttt{clean} command clears the entire working memory. Unlike the implicit
clean-up performed at the beginning of every investigation or control --- which
removes only the proof-tree traces --- this command removes \emph{both} the proof
trees and all the facts asserted by the user during the session. It is therefore
used whenever the expert wants to start a brand-new consultation, with no
information carried over from previous runs:

\begin{lstlisting}
> clean.
\end{lstlisting}


\subsection{Quit}

The \texttt{quit} command closes the user interface and terminates the system.
Before halting, it also clears the working memory, so that no facts or proof
trees are left behind:

\begin{lstlisting}
> quit.
\end{lstlisting}

\chapter{Validation}
\label{chap:validation}

The system is validated with a small regression suite of clinical vignettes,
one per disorder family, collected in the file \texttt{clinical\_tests.pl} and
run with \texttt{swipl clinical\_tests.pl}.

\section{Method}

Each clinical case encodes the answers a clinician would give as a set of
\texttt{known/4} facts. To run the full \emph{investigation} without interactive
input, the harness automatically discovers every symptom the engine may ask ---by
inspecting the clauses of the manifestation, feature and type predicates--- and
sets each of them to \emph{absent} as a baseline; it then overlays the positive
findings of the case, runs the investigation, ranks the disorders by certainty
and checks that the top-ranked diagnosis is the expected one. The
``no-diagnosis'' vignette is considered passed when the highest certainty stays
below $30\%$. Every case is executed under all three uncertainty methods:
\texttt{cf} (minimum t-norm, the default), \texttt{fuzzy} (product t-norm) and
\texttt{fuzzy\_membership} (graded membership functions over duration, frequency
and insight).

\section{Cases and results}

Table~\ref{tab:validation} reports, for each vignette, the certainty assigned to
the expected diagnosis under the three methods (``---'' means the method did not
rank the expected diagnosis first).

\begin{table}[h]
\centering
\caption{Clinical test suite: certainty of the expected diagnosis per method.}
\label{tab:validation}
\small
\begin{tabular}{@{}llccc@{}}
\toprule
\textbf{Vignette} & \textbf{Expected diagnosis} & \textbf{cf} & \textbf{fuzzy} & \textbf{fuzzy-mem.} \\
\midrule
Recurrent unexpected panic attacks & Panic Disorder & 72\% & 57\% & 72\% \\
Assault survivor & Posttraumatic Stress Disorder & 66\% & 28\% & 66\% \\
Contamination obsessions \& rituals & OCD, with poor insight & 73\% & 56\% & 73\% \\
Out-of-proportion spider fear & Specific Phobia, Animal Type & 75\% & 58\% & 75\% \\
Fear of transport / open spaces & Agoraphobia & 75\% & 57\% & 75\% \\
Low mood after job loss & Adjustment Disorder & 70\% & 60\% & 70\% \\
Anxiety during stimulant use & Substance-Induced Anxiety Disorder & 70\% & 37\% & 70\% \\
Spider phobia of $\sim$5 months & Specific Phobia, Animal Type & --- & --- & 64\% \\
Mild non-specific work stress & no strong diagnosis ($<30\%$) & 10\% & 4\% & 10\% \\
\midrule
\textbf{Cases passed} & & \textbf{8/9} & \textbf{8/9} & \textbf{9/9} \\
\bottomrule
\end{tabular}
\end{table}

\section{Findings}

\begin{itemize}
  \item On the seven prototypical cases all three methods rank the correct
  disorder first, and all three correctly refrain from a strong diagnosis on the
  non-specific case.
  \item \emph{Specifiers are surfaced in the result}: the obsessive-compulsive
  case is returned as ``Obsessive-Compulsive Disorder, With Poor Insight''. The
  insight specifier (good or fair / poor / absent or delusional), which applies
  to obsessive-compulsive, body dysmorphic and hoarding disorders, is encoded as
  a selective specifier that labels the diagnosis without changing its presence.
  \item The \texttt{fuzzy} (product) method assigns lower and less stable
  certainties ---most visibly for the polythetic post-traumatic stress disorder
  ($28\%$) and substance-induced anxiety disorder ($37\%$)--- confirming why the
  minimum t-norm is kept as the default.
  \item The borderline vignette is the discriminating case: a marked animal
  phobia present for only about five months. The discrete methods (\texttt{cf}
  and \texttt{fuzzy}) miss the hard ``at least six months'' cut-off and fail to
  reach the diagnosis, whereas \texttt{fuzzy\_membership} grades the five-month
  duration at $0.67$ of the fuzzy set and still ranks the phobia first ($64\%$).
  This illustrates the practical benefit of membership functions at diagnostic
  thresholds, at the cost of requiring a crisp numeric input from the clinician.
\end{itemize}


\begin{thebibliography}{9}

\bibitem{feigenbaum}
E.~A. Feigenbaum,
\emph{The Art of Artificial Intelligence: Themes and Case Studies of Knowledge Engineering},
Proceedings of the 5th International Joint Conference on Artificial Intelligence (IJCAI), 1977.

\bibitem{dsm}
American Psychiatric Association,
\emph{DSM-5-TR. Diagnostic and Statistical Manual of Mental Disorders}, 5th edition, text revision,
American Psychiatric Publishing, 2022.

\bibitem{abc}
\emph{The ABC of Psychopathology. Exploration, identification and treatment of mental disorders}.

\bibitem{stai}
C.~D. Spielberger, R.~L. Gorsuch, R.~E. Lushene,
\emph{Manual for the State--Trait Anxiety Inventory (STAI)},
Consulting Psychologists Press, Palo Alto, CA, 1970.

\end{thebibliography}


\end{document}
